\documentclass{article}
\usepackage{algorithm, algorithmic, amsfonts, amsmath}

\begin{document}
\begin{algorithm}[!h]
\renewcommand{\algorithmicrequire}{\textbf{Input:}}
\renewcommand{\algorithmicensure}{\textbf{Output:}}
	\caption{PU data purification in the wild for object detection}
	\begin{algorithmic}[1]
	    \REQUIRE
	    PU network $N_{pu}$; labeled set $L$; unlabeled set $U$ (massive, from the wild);
	    class prior $\pi_{p}$ (estimable from $L$ and $U$)
	    \ENSURE
	    Teacher detector $D_{t}$; extended dataset $L\cup U^{p}$
	    \STATE \textbf{Step 1: train the PU classifier $N_{pu}$}
	    \WHILE{not converged}
	    \STATE Collect a mini-batch $\{x_{l},x_{u}\}^{N}$ from $L\cup U$
	    \STATE Optimize $N_{pu}$ with the non-negative PU risk
	    $\pi_{p}R_{p}^{+}(g)+\max\!\left(0,\ R_{u}^{-}(g)-\pi_{p}R_{p}^{-}(g)\right)$ % the wild unlabeled pool contains positives, hence nnPU
	    \ENDWHILE
	    \STATE \textbf{Step 2: construct the selected unlabeled dataset}
	    \STATE Use $N_{pu}$ to select images from $U$, forming $U^{p}$ % purification: keep the images that genuinely contain the target class
	    \STATE \textbf{Step 3: generate pseudo labels}
	    \STATE Train a teacher detector $D_{t}$ on the labeled set $L$
	    \STATE Generate pseudo labels on $U^{p}$ using $D_{t}$ % with the noisy images removed, the pseudo-labels are good enough to use
	    \STATE Obtain the extended dataset $L\cup U^{p}$
	    \STATE \textbf{Step 4: self-distillation}
	    \STATE Train the student detector with hint learning and
	    ground-truth-bounded knowledge distillation % the distillation is bounded by ground truth, so the teacher's mistakes are not inherited
		\RETURN $D_{t}$, $L\cup U^{p}$
	\end{algorithmic}
\end{algorithm}
% Guo et al., CVPR 2021, "Positive-Unlabeled Data Purification in the Wild for Object
% Detection". Lines 1 to 11 follow Algorithm 1 of the original; lines 12 and 13 are the self-distillation of Section 3.3.
% Motivation: detection datasets are severely imbalanced. The original gives COCO train2017, where person carries 26k
% boxes and hair drier only 198. Extra data comes from wild unlabeled images, most of which lack the target class,
% so the first step trains a PU classifier to purify rather than a detector.
% Why PU and not ordinary binary classification: the wild unlabeled set contains positives, and treating it wholly as
% negative would introduce bias.
\end{document}
